\section{Core Implementation}
\label{app:code}

This appendix lists the core implementation behind the methods and results of
the thesis: the kernel and bandwidth functions, the KDE chance-constraint
estimators, the residual maps of the three benchmark models, and the
optimization layer that turns these estimators into constraints of a nonlinear
program. Input validation, file input and output, result bookkeeping, and the
code that produces the figures are omitted. Throughout, a residual
$r=g(x,\xi)\leq0$ is safe and $r>0$ is a violation, as in \cref{eq:residual}.
The listings are verbatim excerpts. Their line numbers are those of the source
module named in each caption, so a gap in the numbering marks omitted lines.

\renewcommand{\thelstlisting}{\thesection.\arabic{lstlisting}}

\subsection{KDE Estimators}

\Cref{lst:kernels} evaluates the integrated kernels $F_K$ at a standardized
argument. The Gaussian case uses $F_K=\Phi$, as in \cref{eq:gaussiancase}.
The biased Epanechnikov kernel shifts the integrated Epanechnikov kernel by one
bandwidth, which is the bias $b(h)=h$ in \cref{eq:domination}. Its support
moves to $[0,2]$ in standardized coordinates, so every positive residual
contributes zero safe mass.

\codepart[caption={Integrated Gaussian, Epanechnikov, and biased Epanechnikov
kernels (\texttt{kde\_cco/kde/kernels.py})}, label={lst:kernels}]%
{kde/kernels.py}{54}{88}

\Cref{lst:bandwidth} implements the two bandwidth rules.
\texttt{silverman\_bandwidth} is the robust rule \cref{eq:ch3-bandwidth-rule}
used in the main experiment rows, and \texttt{schuster\_bandwidth} is the scalar
Schuster scaling \cref{eq:schuster-local-bandwidth} used in the paper-aligned
profile. Both keep a positive floor, so the kernel never becomes singular inside
the optimizer.

\codepart[caption={Silverman and Schuster bandwidth rules
(\texttt{kde\_cco/kde/bandwidth.py})}, label={lst:bandwidth}]%
{kde/bandwidth.py}{8}{62}

\Cref{lst:scalar} turns residual samples $r_j=g(x,\xi_j)$ into the quantities
the optimizer consumes. \texttt{kde\_safe\_probability} is $\hat p_h(x)$ of
\cref{eq:integrated}, imposed as \cref{eq:smoothcc}.
\texttt{kde\_violation\_upper} with \texttt{bias=True} returns
$1-\hat p_h^{\,b}(x)$ of \cref{eq:biasedcc}, which is never smaller than the
empirical violation frequency. \texttt{kde\_safe\_probability\_gradient} is
the analytic gradient \cref{eq:gradient} for the Gaussian kernel.

\codepart[caption={Integrated-KDE chance-constraint estimators and gradient
(\texttt{kde\_cco/constraints/scalar.py})}, label={lst:scalar}]%
{constraints/scalar.py}{86}{149}

\Cref{lst:joint} reduces the residual vector $(g_1,\dots,g_m)$ to one scalar
statistic. \texttt{max\_violation} is the exact reduction
\cref{eq:exactreduction}, and \texttt{smooth\_max} is the log-sum-exp surrogate
\cref{eq:lse} with unit weights $w_r=1$ and temperature $\tau$, evaluated stably
through \texttt{scipy.special.logsumexp}.

\codepart[caption={Exact and smoothed joint statistics
(\texttt{kde\_cco/constraints/joint.py})}, label={lst:joint}]%
{constraints/joint.py}{26}{46}

\subsection{Benchmark Models}

\Cref{lst:static-model} defines the static nonlinear benchmark of
\cref{ch:experiments} (\cref{tab:local-static-results}). The objective is a
shifted quadratic with a small coupling term, and the deterministic constraints
keep the decision in a compact nonnegative region. The uncertain residual is
the circular safety condition
$(x_1+0.25\,\xi_1)^2+(x_2+0.25\,\xi_2)^2-1\leq0$, evaluated for all samples at
once.

\codepart[caption={Static nonlinear benchmark
(\texttt{kde\_cco/problems/static\_problem.py})}, label={lst:static-model}]%
{problems/static_problem.py}{40}{63}

\Cref{lst:dispatch-model} defines the two-generator dispatch benchmark
(\cref{tab:dispatch-results}). The cost is quadratic in the output of each
generator. The first residual component is the balance deficit
$d-r_0-\xi-(x_1+x_2)$, where $d$ is the demand, $r_0$ the nominal renewable
output, and $\xi$ the renewable forecast error. The two capacity components do
not depend on $\xi$, so the stochastic part of the problem is the balance
equation alone.

\codepart[caption={Dispatch benchmark
(\texttt{kde\_cco/problems/energy\_dispatch.py})}, label={lst:dispatch-model}]%
{problems/energy_dispatch.py}{53}{80}

\Cref{lst:lunar-model} gives the lunar model of \cref{sec:lunar-benchmark}.
\texttt{dynamics} is \cref{eq:lunar-dynamics}, and \texttt{terminal\_residual}
and \texttt{path\_residual\_nodes} are the residuals inside the terminal and
path events \cref{eq:lunar-terminal-risk,eq:lunar-path-risk}. When all $n_t$
node controls are passed, \texttt{objective} evaluates the trapezoidal fuel
integral whose mismatch with the Euler defects is \cref{eq:quadrature-mismatch}.

\codepart[caption={Lunar landing dynamics, objective, and residuals
(\texttt{kde\_cco/problems/lunar\_landing.py})}, label={lst:lunar-model}]%
{problems/lunar_landing.py}{78}{107}
\codepart[nolol]{problems/lunar_landing.py}{117}{123}
\codepart[nolol]{problems/lunar_landing.py}{139}{145}

\subsection{Optimization}

\Cref{lst:static-solver} connects a chance constraint to SciPy's SLSQP. A
chance constraint is passed as a function that is nonpositive when it holds,
for example $(1-\varepsilon)-\hat p_h(x)$. SciPy instead expects inequality
functions $c(x)\geq0$, so \texttt{\_constraint\_from\_chance} negates the
function, and \texttt{\_solve} passes it together with the deterministic
constraints and the box bounds to \texttt{scipy.optimize.minimize}.

\codepart[caption={Chance-constraint adapter for SLSQP
(\texttt{kde\_cco/solvers/static.py})}, label={lst:static-solver}]%
{solvers/static.py}{107}{127}
\codepart[nolol]{solvers/static.py}{152}{164}

\Cref{lst:collocation} is the optimal-control transcription of
\cref{ch:optimalcontrol}. \texttt{transcribe\_euler} returns the forward Euler
defects \cref{eq:euler-defect} and the boundary residuals. The inner function
\texttt{values} of \texttt{\_chance\_constraints} computes the terminal and
node-wise path margins of one continuation stage. The nominal stage uses the
sample mean, the worst-case stage the largest residual, and the scenario stage
one inequality per sample. The unbiased stage uses the Gaussian KDE and the
biased stage the shifted Epanechnikov kernel, both in the safe-probability form
equivalent to \cref{eq:ocp-kde}. The core of \texttt{solve\_collocation}
stacks states and controls into one decision vector and passes the defects as
equalities and the chance margins as inequalities to SLSQP.

\codepart[caption={Euler transcription and KDE chance stages
(\texttt{kde\_cco/solvers/collocation.py})}, label={lst:collocation}]%
{solvers/collocation.py}{17}{39}
\codepart[nolol]{solvers/collocation.py}{77}{104}
\codepart[nolol]{solvers/collocation.py}{144}{173}

\Cref{lst:continuation} is the bandwidth continuation of the paper-aligned
profile (\cref{tab:paper-aligned-sequence}). \texttt{\_continuation} solves a
sequence of problems with bandwidths $c\,h_N$, warm-starts each solve from the
previous decision, and stops at the first failure. The paper-aligned unbiased
KDE uses one Gaussian solve at the bandwidth $h_N$ of
\cref{eq:schuster-local-bandwidth}. The paper-aligned biased KDE runs the
multipliers $c\in\{0.25,0.5,1,2\}$ with the $b(h)=h$ Epanechnikov surrogate and
keeps the last successful stage.

\codepart[caption={Bandwidth continuation in the paper-aligned profile
(\texttt{kde\_cco/experiments/paper\_aligned.py})}, label={lst:continuation}]%
{experiments/paper_aligned.py}{41}{53}
\codepart[nolol]{experiments/paper_aligned.py}{73}{100}
